\documentclass{article}

\usepackage[utf8]{inputenc}
\usepackage[margin=1in]{geometry}
\usepackage{booktabs}
\usepackage[hidelinks]{hyperref}

\title{Premise-Grounded Detection of Unfaithful\\
Chain-of-Thought Reasoning}
\author{Suramya Raj Angdembay \and Dikshant Aryal}
\date{Summer 2026}

\begin{document}

\maketitle

\section{Research Context and Problem Statement}

Large language models (LLMs) increasingly solve hard problems by emitting a
chain-of-thought (CoT): a step-by-step natural-language explanation that precedes the
final answer. These explanations are widely treated as a window into how a model
reasons, and they are becoming central to AI-safety practices such as CoT monitoring.
A fluent CoT can nonetheless be \emph{unfaithful}: the stated reasoning may not
actually support, or be responsible for, the answer the model gives. A model may reach
an answer before constructing its explanation, rely on an unstated or invalid premise,
ignore evidence it appears to cite, or arrive at a correct answer through an invalid
path~\cite{lanham2023measuring, cox2026decoding}. Because final-answer accuracy cannot
reveal any of these failures, accuracy alone is an insufficient measure of trustworthy
reasoning.

This matters wherever reasoning is used to justify a decision---scientific, medical,
legal, or educational---and wherever CoT is relied upon for oversight: if explanations
cannot be trusted to reflect the reasoning that produced an answer, they cannot be used
to audit, debug, or safely supervise a model. The recent FaithCoT-Bench benchmark shows
that detecting instance-level CoT unfaithfulness remains difficult, and that current
detectors are weakest precisely where they matter most---on knowledge-intensive
questions and on the convincing traces produced by stronger
models~\cite{shen2026faithcot}.

\textbf{Problem.} Given a question, a generated CoT, and a final answer, we will
(1)~predict whether the reasoning trace is faithful, (2)~localize the first step that
becomes unsupported or causally irrelevant, and (3)~identify the premise or evidence
item responsible for the failure. We target \emph{observable trace faithfulness}---%
whether the stated reasoning is sufficiently supported and affects the answer as
claimed---rather than the stronger mechanistic claim that the text exactly describes the
model's internal computation. The specific gap we address is that existing detectors
judge a CoT holistically and perturb it without regard to which parts of the reasoning
are actually load-bearing.

\section{Proposed Solution}

We propose a \emph{premise-grounded} faithfulness detector that converts a free-form CoT
into an auditable dependency structure and verifies it locally. The pipeline has five
stages:

\begin{enumerate}
  \item \textbf{Decompose} the CoT into atomic reasoning steps, each expressing a single
        claim, inference, or calculation.
  \item \textbf{Build a premise graph.} Following Premise-Augmented Reasoning Chains
        (PARC), we represent the trace as a directed acyclic graph in which each step
        points to the earlier steps, given facts, and external evidence it depends
        on~\cite{mukherjee2025parc}.
  \item \textbf{Verify local support} for each step against \emph{only} its declared
        premises, using two complementary channels: a \emph{logical} channel (a
        premise-restricted LLM judge, with symbolic execution where the step is
        formalizable) and a \emph{factual} channel (natural-language inference and
        retrieved external evidence, behind a retrieval-quality gate that separates
        retrieval failure from reasoning failure)~\cite{asai2024selfrag}.
  \item \textbf{Test causal relevance} with graph-guided counterfactual interventions:
        we use the graph to select load-bearing premises, then remove, replace, or
        contradict them and observe whether the downstream reasoning and answer change
        appropriately. Because the targets are chosen from the model's own declared
        dependencies, these interventions should yield a stronger faithfulness signal
        than the random or position-based perturbations used by prior
        work~\cite{lanham2023measuring}.
  \item \textbf{Aggregate} the structural, evidential, and interventional signals into an
        instance-level faithfulness label plus a first-unsupported-step localization, via
        a transparent weighted score or lightweight classifier (no frontier-scale
        training).
\end{enumerate}

This design \emph{extends}, rather than re-derives, recent structured-verification work.
PARC and the Graph-of-Verification framework already build premise/dependency graphs and
verify steps under their premises, but for \emph{error correctness} in mathematics and
with an LLM judge alone~\cite{mukherjee2025parc, fang2025gov}; VeriCoT adds neuro-symbolic
per-step premise checking but likewise targets logical validity rather than
faithfulness~\cite{feng2025vericot}. Our contribution is to (a)~repurpose the premise
graph to \emph{select} counterfactual interventions on load-bearing premises, (b)~ground
step verification in retrieved evidence for the knowledge-intensive regime where holistic
judges fail~\cite{shen2026faithcot}, and (c)~unify instance-level detection with
first-unsupported-step localization. The method is model-agnostic and
black-box/gray-box; a white-box validation using linear probes is an optional later
extension.

\section{Evaluation and Implementation Plan}

\subsection{Evaluation Plan}

\textbf{Datasets.} Our primary benchmark is FaithCoT-Bench/FINE-CoT, the only
instance-level CoT-faithfulness benchmark with expert annotations and step-level
evidence, spanning logical (LogiQA), factual (TruthfulQA), mathematical (AQuA), and
knowledge-intensive biomedical (HLE-Bio) reasoning~\cite{shen2026faithcot}. We will test
cross-dataset generalization on the parallel step-level faithfulness benchmark
GRACE~\cite{pham2026grace}, evaluate step-level localization against process-error
benchmarks (ProcessBench, REVEAL)~\cite{zheng2024processbench, jacovi2024reveal}, and use
datasets with machine-verified premises (FOLIO) to obtain gold premise graphs and clean
premise-removal counterfactuals~\cite{han2024folio}. Where faithfulness labels are
unavailable, results on synthetic-corruption and human-labeled data will be reported
separately, since synthetic errors are systematically easier to detect.

\textbf{Conditions and baselines.} We will compare the full detector against:
final-answer correctness; whole-trace and step-level LLM-as-judge; untargeted
(random/position-based) perturbation~\cite{lanham2023measuring}; and premise-graph
verification without interventions or evidence (a PARC/GoV-style
baseline)~\cite{mukherjee2025parc, fang2025gov}. Ablations will remove graph-guided
intervention, remove evidence retrieval, substitute gold premise links for inferred
ones, vary the premise extractor's size, and use a cross-model judge to control for
self-preference bias.

\textbf{Metrics.} \emph{Detection:} macro-F1, precision/recall for unfaithful traces,
AUROC/AUPRC, Cohen's~$\kappa$, and calibration error. \emph{Localization:} first-error
accuracy, step-level precision/recall/F1, and premise-link and evidence-attribution
accuracy. \emph{Intervention:} answer-flip rate when a load-bearing premise is perturbed,
(in)appropriate stability under critical versus irrelevant changes, and counterfactual
consistency measured as a distributional answer shift. \emph{Cost:} tokens,
verifier/retrieval calls, latency, and estimated cost per example.

\textbf{Success criteria.} The project succeeds if the premise-grounded detector achieves
at least one of: (i)~a $\geq 5$-point absolute macro-F1 gain over a strong judge baseline
on a knowledge-intensive subset; (ii)~a clear improvement in first-error or step-level
localization; (iii)~substantially better sensitivity to critical perturbations without
increased sensitivity to irrelevant ones; or (iv)~comparable detection using a smaller,
cheaper judge. We will report confidence intervals and paired significance tests and---%
following critiques that faithfulness metrics can track answer-choice bias---control for
accuracy and option bias throughout. Small gains in answer accuracy alone will not be
treated as success.

\subsection{Timeline}

We organize the work as a roughly twelve-week summer effort
(Table~\ref{tab:timeline}). Each phase produces a concrete, evaluable deliverable, so
that a useful result exists even if later phases are cut for time. The minimum viable
study runs through Phase~3; the white-box probe in Phase~4 is a stretch goal.

\begin{table}[h]
    \centering
    \caption{Proposed summer timeline (approximately twelve weeks).}
    \label{tab:timeline}
    \begin{tabular}{p{0.17\textwidth} p{0.17\textwidth} p{0.54\textwidth}}
        \toprule
        \textbf{Phase} & \textbf{Period} & \textbf{Milestones / deliverables} \\
        \midrule
        0.~Setup \& baselines & Weeks 1--2 (late June) &
        Reproduce FaithCoT-Bench baselines; stand up the PARC code and PERL data;
        identify the weakest task/model combinations. \\
        1.~Structural detector & Weeks 3--5 (July) &
        CoT segmentation; premise-link prediction; premise-restricted two-channel local
        verification (LLM judge + NLI). \\
        2.~Evidence \& interventions & Weeks 6--8 (late July--Aug.) &
        Evidence retrieval with a relevance gate for factual premises; graph-guided
        counterfactual interventions; symbolic checking on the logic (FOLIO) subset. \\
        3.~Evaluation \& ablations & Weeks 9--10 (mid-Aug.) &
        Full comparison against baselines; ablations; paired significance tests; cost and
        qualitative error analysis. \\
        4.~Writing \& stretch & Weeks 11--12 (late Aug.--Sept.) &
        Workshop-style write-up; optional white-box (linear-probe) validation if time
        permits. \\
        \bottomrule
    \end{tabular}
\end{table}

\section{Team}

\begin{table}[h]
    \centering
    \caption{Research team}
    \label{tab:team}
    \begin{tabular}{ll}
        \toprule
        \textbf{Name} & \textbf{Role} \\
        \midrule
        Suramya Raj Angdembay & Undergraduate researcher \\
        Dikshant Aryal & Undergraduate researcher \\
        Dr. Nick Rahimi & Faculty advisor / graduate mentor \\
        \bottomrule
    \end{tabular}
\end{table}

\bibliographystyle{plain}
\bibliography{ur2phd}

\end{document}
